% Não altere nada nesta parte
\begin{center}
  {\bfseries \sffamily \fontsize{11}{11} \TítuloDoTrabalho}
\end{center}

\noindent{\bfseries \sffamily\large Resumo}
\addcontentsline{toc}{chapter}{Resumo}


O trabalho tem como objetivo investigar a aplicação conjunta 
de modelos de Large Language Models (LLMs) e Knowledge Graphs na 
análise de comunicações do Federal Reserve, com foco na identificação 
de sinais relacionados à orientação da política monetária. Foi utilizado 
como fontes documentos oficiais do Federal Reserve, por exemplo comunicados 
do FOMC, atas de reuniões, discursos, testemunhos e press releases. 
Inicialmente, foram coletados 1.283 links de comunicações públicas a partir 
do site oficial da instituição, seguidos pela extração e preparação dos textos 
em formato estruturado. Em seguida, modelos de linguagem foram utilizados para 
extrair sinais monetários e classificá-los em categorias de postura monetária, 
variando entre dovish e hawkish. Cada sinal extraído também recebeu uma 
pontuação numérica e um vetor de embedding, permitindo o cálculo de 
similaridade semântica. A partir desses dados, foi construída uma versão 
inicial de um knowledge graph, composto por publicações, datas, sinais 
monetários e relações semânticas entre sinais. Os resultados preliminares 
indicam a viabilidade metodológica da abordagem proposta, ao integrar 
processamento de linguagem natural, embeddings, grafos de conhecimento 
e algoritmos de análise de redes para apoiar a compreensão da comunicação 
monetária ao longo do tempo.

\vspace{1cm}
\noindent\textbf{Palavras-chave:} 
Federal Reserve; política monetária; modelos de linguagem; 
grafos de conhecimento; processamento de linguagem natural.
